\section{Auditoría de fuentes: lo que esta clase trata en realidad es «cómo se evalúa la distribución generada»}

La conferencia de Mark Chen va desde el 3-gram hasta GPT-3, iGPT, DALL-E y Codex. En apariencia es una cronología de modelos, pero el hilo conductor real es el cambio constante de la evaluation interface. En la etapa de texto se comparan la sample coherence y la capacidad humana de reconocimiento. En la etapa de aprendizaje no supervisado se compara la transfer. En la etapa de imágenes se observan a la vez la completion y la representation. En la etapa de código, los resultados generados deben pasar unit tests y se usa pass@k para evaluar múltiples muestreos.

El video original es una slide talk a pantalla completa y el repositorio no tiene un deck independiente. Este documento recupera 45 teaching slides a partir de las transiciones del video oficial. El bumper de Stanford, tres páginas que solo son separadores, los agradecimientos y el thank-you repetido se marcan como optional. La clase se impartió en 2021. Las descripciones de productos y capacidades de Codex, GPT-3 y DALL-E se tratan según ese momento histórico y no deben reescribirse como si describieran el estado actual de los productos en 2026.

\begin{importantbox}{Los tres ciclos cerrados de esta clase}
\begin{enumerate}
  \item \textbf{Generation loop}: aprender la distribución $\rightarrow$ muestreo $\rightarrow$ evaluación humana/de la tarea $\rightarrow$ ajuste de la scale y la recipe;
  \item \textbf{Transfer loop}: unlabeled data $\rightarrow$ autoregressive objective $\rightarrow$ representation/behavior $\rightarrow$ downstream task;
  \item \textbf{Code loop}: problem prompt $\rightarrow$ sampled programs $\rightarrow$ unit tests $\rightarrow$ pass@k $\rightarrow$ temperature/model decision.
\end{enumerate}
\end{importantbox}

\begin{warningbox}{Sample quality, benchmark score y deployment reliability son cosas distintas}
Que un texto parezca una noticia no prueba que sea verdadero. Que una imagen corresponda al prompt no prueba que haya generalización composicional. Que un programa pase los benchmark tests no prueba que sea seguro o mantenible, ni que cumpla requisitos ocultos. Este documento distingue en todo momento entre «en la distribución del modelo existe una muestra correcta», «la decoding actual encontró esa muestra» y «el sistema desplegado puede confiar en la salida».
\end{warningbox}

\begin{knowledgebox}{Términos clave: los cuatro objetos de un sistema generativo}
\textbf{Model distribution} es la distribución de probabilidad que el modelo asigna a las salidas candidatas. \textbf{Sample} es una salida concreta obtenida de esa distribución. \textbf{Verifier} juzga a los candidatos mediante reglas, tests o evaluación humana. \textbf{Selector} elige la salida final entre varios candidatos. El modelo, el muestreo, la verificación y la selección pueden mejorarse por separado, así que el éxito final no puede atribuirse por completo a la escala de parámetros.
\end{knowledgebox}

\subsection{Resumen de la sección}

Esta no es solo una clase sobre la historia de GPT. Es una clase sobre cómo se evalúa el modelado generativo. A medida que la salida pasa del texto a las imágenes y al código, la evaluación tiene que dejar de medir la similitud superficial y pasar a medir la evidencia de la tarea y la distribución de muestreo.
